\scanpage{34}
\uline{Prop}
\[
\left\{\begin{aligned}
P^i(X)&\xrightarrow{\sim}P^i(Y) &&\text{si }i\leq n-2\\
P^{n-1}(X)&\hookrightarrow P^{n-1}(Y) &&\deleted{\unsure{\text{si }i}}\\
P^i(X)&=\operatorname{Ker}\varphi^i &&\text{si }i\geq n
\end{aligned}\right.
\]
\diagram{\begin{tikzcd}[ampersand replacement=\&,column sep=large,row sep=large]
\deleted{\unsure{P}}H^i(X)\arrow[r,"L_X^{n-i+1}"]\arrow[d,"\varphi^i"',"\unsure{\wr|2}"] \& H^{2n-i+2}(X)\\
H^i(Y)\arrow[r,"L_Y^{n-i\deleted{\unsure{+1}}}"'] \& H^{2(n-1)\unsure{-|+}i}(Y)\arrow[u,"\varphi_{2(n-1)+i}"',"\unsure{\wr|2}"]
\end{tikzcd}}
\[\left.\vphantom{\begin{matrix}H^i(X)\\H^i(Y)\end{matrix}}\right]\quad\deleted{\unsure{\forall}}\ i\leq n-2\]
\[i\leq n-2\Rightarrow 2(n-1)-i\geq n.\]
\diagram{\begin{tikzcd}[ampersand replacement=\&,column sep=large,row sep=large]
H^j(X)\arrow[d,"\varphi^*"']\arrow[dr,"L_X"] \& \\
H^j(Y)\arrow[r,"\sim"]\&H^{j+2}(X)
\end{tikzcd}}
\[j\geq n\]

\uline{\unsure{Cor}}\qquad Notons
\[
\left\{\begin{aligned}
H^*(X)&=\deleted{\unsure{\sum_i}}\coprod_{\substack{i+j\leq n\\i,j\geq0}}L_X^jP^i(X)
=\coprod_{\substack{i+j\leq n-1\\i,j\geq0}}L_X^jP^i(X)\oplus\coprod_{j\geq n}P^j(X)\\
H^*(Y)&=\coprod_{\substack{i+j\leq n-1\\i,j\geq0}}L_Y^jP^i(Y).
\end{aligned}\right.
\]
\uline{Cor}\qquad $\varphi^*$ induit un mor d'objets bigradués
\[
\deleted{\unsure{=}}\coprod_{\substack{i+j\leq n-1\\i,j\geq0}}L_X^jP^i(X)
\longrightarrow H^*(Y)=\unsure{\coprod}_{\substack{i+j\leq n-1\\i,j\geq0}}L_Y^jP^i(Y);\]
\deleted{\unsure{qui est un}}\inserted{\unsure{c'est}} \unsure{induit} un isom sur les morceaux bihomogènes, sauf
\[
L_X^0P^{n-1}(X)=P^{n-1}(X)\longrightarrow\deleted{\unsure{P^{n-1}}}L_Y^0P^{n-1}(Y)=P^{n-1}(Y).
\]
\uline{Cor}\qquad Le noyau de $\varphi^*$ est \inserted{\unsure{égal à celui de $L_X$, i.e.}}
\[\deleted{\unsure{\sum}}\coprod_{j\geq n}P^i(X)\]
\deleted{(qui est \unsure{un idéal})}. Le conoyau est
\[
\begin{aligned}
E^{n-1}&=H^{n-1}(Y)/\operatorname{Im}H^{n-1}(X)=P^{n-1}(Y)/\operatorname{Im}P^{n-1}(X)\\
&\simeq\operatorname{Ker}\bigl(\varphi_{n-1}:H^{n-1}(Y)\mathrel{\text{---}}H^{n+1}(X)\bigr)
\end{aligned}
\]
\mspage{33}
